\begin{exercício}{Convergência da sequência de raio espectral}{ex6}
Seja \(X\) um espaço de Banach complexo, \(T \in \bounded(X)\) um operador linear limitado e seja \(\family{a_n}{n\in \mathbb{N}} \subset \mathbb{R}_+\) uma sequência tal que \(a_{n+m} \leq a_n + a_m\) para todos \(n,m \in \mathbb{N}.\)
\begin{enumerate}[label=(\alph*)]
  \item Mostre que \(\lim_{n \to \infty}{\frac{a_n}{n}}= \inf_{n\in \mathbb{N}}{\frac{a_n}{n}}.\)
  \item Mostre que a sequência \(\family{\norm{T^n}^{\frac1n}}{n\in\mathbb{N}}\) converge para o seu ínfimo.
\end{enumerate}
\end{exercício}
\begin{proof}[Resolução de (a)]
  Seja \(b : \mathbb{N} \to \mathbb{R}_+\) a sequência definida por \(n \mapsto \frac{a_n}{n},\) então é claro que essa sequência é limitada inferiormente por zero, portanto seu ínfimo \(L\) é não negativo e finito. Seja \(\epsilon > 0,\) então existe \(m \in \mathbb{N}\) tal que
  \begin{equation*}
      b_m \leq L + \epsilon.
  \end{equation*}
  Seja \(n \geq m,\) então existem \(q \in \mathbb{N}\) e \(r \in \mathbb{N}\) com \(r < m\) tais que \(n = qm + r,\) portanto pela subaditividade da sequência \(a_n\), segue que
  \begin{equation*}
    b_n = \frac{a_n}{n} \leq \frac{q a_m}{n} + \frac{a_r}{n} \leq \frac{q a_m}{qm} + \frac{a_r}{n} = b_m + \frac{a_r}{n}.
  \end{equation*}
  Seja \(\eta > 0\) e escrevamos \(M = \max\setc{a_k}{k \in \mathbb{N} : k < m},\) portanto
  \begin{equation*}
    n > \max\set{\eta^{-1} M, m} \implies b_n < b_m + \eta,
  \end{equation*}
  isto é,
  \begin{equation*}
    \limsup_{n\to\infty}{b_n} \leq b_m.
  \end{equation*}
  Como \(\epsilon\) é arbitrário, temos então
  \begin{equation*}
    \limsup_{n\to\infty}{b_n} \leq \inf_{m \in \mathbb{N}}{b_m} = L.
  \end{equation*}
  Agora, como
  \begin{equation*}
    L = \inf_{n \in \mathbb{N}}{b_n} \leq \liminf_{n\to\infty}{b_n} \leq \limsup_{n\to\infty}{b_n} \leq L,
  \end{equation*}
  sabemos que \(b\) é convergente e \(b_n \to L\).
\end{proof}
\begin{proof}[Resolução de (b)]
  Seja \(t : \mathbb{N} \to \mathbb{R}_+\) a sequência \(n \mapsto \norm{T^n}.\) Notemos que
  \begin{equation*}
    \norm{T^{n + m}} \leq \norm{T^n} \norm{T^m} \implies t_{n+m} \leq t_n t_m
  \end{equation*}
  para todos \(n, m \in \mathbb{N}.\) Podemos sem perda de generalidade assumir que \(T\) não é nilpotente, já que se fosse o caso, é claro que a sequência \(t\) converge para o seu ínfimo, a saber, zero. Consideramos a sequência \(a : \mathbb{N} \to \mathbb{R}_+\) dada por \(n \mapsto \ln t_n,\) que satisfaz a condição de subaditividade,
  \begin{equation*}
    a_{n + m} = \ln(t_{n+m}) \leq \ln(t_n t_m) = \ln t_n + \ln t_m = a_n + a_m,
  \end{equation*}
  portanto temos
  \begin{equation*}
    \frac{a_n}{n} \to \inf_{n \in \mathbb{N}}{\frac{a_n}{n}} \implies \ln\left(\norm{T^n}^{\frac1n}\right) \to \inf_{n\in \mathbb{N}}{\ln\left(\norm{T^n}^{\frac1n}\right)}
  \end{equation*}
  Da continuidade da exponencial e de sua monotonicidade crescente, segue então que
  \begin{equation*}
    \lim_{n\to \infty}{\norm{T^n}^{\frac{1}{n}}} = \inf_{n \in \mathbb{N}}{\norm{T^n}^{\frac1n}}
  \end{equation*}
  como desejado.
\end{proof}
